\section{Scope, status, and principal results}
\label{sec:introduction}

The ProveIt Polylogarithms manuscript~\cite{ProveIt} organizes a large collection of
functional equations, special values, experimental reductions, and
arithmetic bridges. This continuation concentrates on a connected part
of that programme: logarithmic gamma moments, Herglotz derivatives, and
depth-two cyclotomic polylogarithms. The common methodological question is
how to turn a useful formal representation into a statement with
controlled errors or a precisely specified algebraic meaning.

The source baseline is commit
\begin{center}\small
\texttt{78f9e5df05eb0abb16c89fb00b9e89403598bd17}
\end{center}
of \nolinkurl{VladimirReshetnikov/ProveIt}. We concentrate on Chapters~4, 7, 9,
and~10 in \nolinkurl{Analysis/Polylogarithms/docs/manuscript}.
We also inspected the existing October continuations on alternating
harmonic sums, corpus corrections, rational-grid distribution ranks,
Stieltjes derivatives, and Herglotz obstructions. The accompanying source
manifest records the examined files. This article proposes additions and
corrections; it does not modify the upstream repository.

\subsection{What is proved here}

There are three principal developments.

First, the classical moment expansion
\[
 \frac1{n!}\int_0^1\log^n\Gamma(x)\dd x
 \sim\sum_{k\geq1}\frac{(-1)^{k-1}b_k}{k^n}
\]
has coefficients that grow exponentially. We identify their precise
growth and prove a sharp error law when the number of retained terms
grows linearly with \(n\). At the least-term scale the actual error is
asymptotic to \emph{one half of the first omitted term}; the article
computes two correction terms. This is stronger than a fixed-length
asymptotic expansion or a heuristic minimum-term estimate. The proof
requires a local analytic continuation of an inverse gamma branch and a
uniform estimate for a Taylor remainder across its circle of convergence.

Second, for
\[
 F(x)=\sum_{n\geq1}\frac{\psi(nx)-\log(nx)}n,\qquad
 D_r(x)=\frac{(-1)^{r+1}x^{r+1}}{r!}F^{(r)}(x),
\]
we show that
\[
 D_r(x)=\zeta(2)+\frac xr(\log x-H_{r-1})+E_r(x).
\]
The displayed elementary expression contains every algebraic order of
the large-\(r\) expansion. The remaining term admits a divisor-weighted
Tricomi-\(U\) representation, an effective gamma-factor bound, and an
explicit oscillatory asymptotic of size
\(r^{-3/4}e^{-2\sqrt{\pi r/x}}\). This yields the sharp exponential rate
and infinitely many positive and negative remainders. Here the derivative
order tends to infinity with \(x\) fixed; this is a different limit from
the familiar large-\(x\) Bernoulli expansion.

Third, we determine all-weight dimensions of a \emph{specified formal
presentation} of imaginary depth-two values at levels 3 and 4, after
quotienting by single products and the manuscript's \(g\)-family.
The Hilbert series are
\[
 \frac{t^5}{(1-t^2)(1-t^6)}
 \quad\hbox{and}\quad
 \frac{t^3}{(1-t^2)^2},
\]
respectively. A sparse functional separates the proposed Gaussian
\(S_4\) reduction from every row in that presentation. This identifies
what the chosen relation family cannot establish. It is not a proof of
linear independence of actual periods, and it is not a disproof of the
conjecture.

These results are supported by exact rational row calculations,
independent integral evaluations, precision increases, and explicit
analytic tail bounds. Their proofs do not depend on the numerical
checks.

\subsection{Classical input and attribution}

The moment expansion itself is due to Amdeberhan, Coffey, Espinosa,
Koutschan, Manna, and Moll~\cite{ACEMKM}. The square-root mechanism used
for its late coefficients is classical analytic combinatorics
\cite{FS}; its hypotheses and the constants are established here for
the inverse relevant to the moments. The growing-index remainder
analysis is the additional step.

Radchenko and Zagier~\cite{RZ} already state closure of all rational
Herglotz derivatives in depth-one cyclotomic polylogarithms. We supply
an explicit finite Bernoulli--Hurwitz coefficient formula and a
regularization calculation, not a new claim of closure. Their arithmetic
integral representations also underlie the divisor expansion used here.
Large-parameter asymptotics of Tricomi's function are classical
\cite{DLMF,Temme}; the effective sum over arithmetic sectors and its
high-derivative consequences are developed in this article.

The cubic log-gamma moment already has a Tornheim--Witten derivative
evaluation in Bailey--Borwein--Borwein~\cite{BBB}. We correct the
manuscript's open-problem wording and independently derive an equivalent
formula, including convergence details and an evaluator. A reduction to
a more restrictive prescribed basket of ordinary zeta derivatives
remains a different question.

The all-depth odd-harmonic-sum reflection and several low-weight
Gaussian reductions have already been proved in the repository's
continuations. We use them as established context. Likewise, a count of
survivors in a relation matrix is always tied to the exact relations
imposed. Zhao's broader work on cyclotomic relations~\cite{Zhao}
provides a useful warning against identifying a restricted row space
with the full algebra of periods.

\subsection{Conventions and verification levels}

Throughout, \(\log\) is the real logarithm on positive real arguments;
complex powers and the function \(U\) use principal branches unless
otherwise specified. We write \(H_n=\sum_{j=1}^n j^{-1}\), \(H_0=0\),
\(\psi=\Gamma'/\Gamma\), and
\[
 \Li_{a,b}(u,v)=\sum_{m>n\geq1}\frac{u^m v^n}{m^a n^b}.
\]
At convergent boundary points this is the ordinary value. When a
leading \(\Li_1(1)\) is required in a relation, the regularization is
specified explicitly. The Dirichlet beta and eta functions are
\[
 \beta(s)=\sum_{n\geq0}\frac{(-1)^n}{(2n+1)^s},\qquad
 \eta(s)=\sum_{n\geq1}\frac{(-1)^{n-1}}{n^s},
\]
with \(\eta(1)=\log2\) and \(G=\beta(2)\).

A theorem in this article has an analytic or exact algebraic proof.
A formal dimension is the dimension of the defined presentation.
An integer-relation candidate is a conjecture even when its residual
is extraordinarily small. Most numerical diagnostics use ordinary
high-precision arithmetic, so analytic truncation estimates are
distinguished from machine-certified rounding enclosures. The
\(S_6\) residual certificate instead uses exact rational arithmetic
with proved tails. No claim
of a formal proof-assistant verification is made.

The novelty assessment is relative to the pinned manuscript, the
inspected continuations, and the cited primary literature. The article
does not claim an exhaustive priority search. All statements presented
as theorems stand on their supplied proofs, independently of that
bibliographic qualification.

